\chapter*{Abstract}
\addcontentsline{toc}{chapter}{Abstract}

Large language models (LLMs) can communicate a misleading financial risk impression without making an explicit false statement. A response may omit a material adverse fact, weaken its specificity, bury it after benefits, or replace it with a generic caveat. Such behaviour is poorly captured by evaluations centred on hallucination, instructed lying, or agentic scheming under an explicit goal conflict. This thesis proposes \benchmarkname, a finance-domain benchmark for measuring \emph{risk concealment} under deployment-realistic, benign prompting.

The planned benchmark contains \placeholder{40 scenario families, subject to final review}, split between single-turn and multi-turn interactions. Each scenario is evaluated in a $3\times3$ design crossing neutral, production-baseline, and production-integrity prompts with neutral, anxious risk-averse, and positive risk-seeking user personas. Target models receive a realistic source artefact and an annotated pool of atomic facts. Outputs are scored separately for false claims, selective omission, misleading framing, risk de-emphasis, and downstream distortion in a simulated user's beliefs or actions. The construct is behavioural: it measures what happens to available risk information, not whether a model possesses a human-like intention to deceive.

The proposed analysis combines deterministic fact-level measures, constrained LLM-assisted annotation, human validation, and mixed-effects models accounting for repeated observations within scenarios and families. \placeholder{Insert the preregistered primary endpoint, model roster, run configuration, power analysis, and human-audit sample before data collection.} \placeholder{After experiments, insert numerical findings on benign-condition prevalence, dominant failure mode, integrity-prompt effects, persona interactions, model heterogeneity, annotation reliability, and simulated user harm, with confidence intervals.}

The intended contribution is to shift evaluation from ``did the model say something false?'' to ``did the model enable the user to understand the material risk?''
